\documentclass[11pt,a4paper]{article}
\usepackage[utf8]{inputenc}
\usepackage{amsmath}
\usepackage{amsfonts}
\usepackage{amssymb}
\usepackage{amsthm}
\usepackage[margin=2cm]{geometry}
\usepackage{enumitem}

\begin{document}

\begin{center}
  \Large\textbf{Soal Pembinaan OSITS 2026}
\end{center}
\vspace{0.5cm}

\section*{Analisis Real}
\begin{enumerate}[label=\arabic*.]
  \item % 1.6.29
        Buktikan bahwa
        \[
          \lim_{n \to \infty} \left(1+\frac{1}{n}\right) \left(1+\frac{2}{n}\right)^{1/2} \dots \left(1+\frac{n}{n}\right)^{1/n} = e^{\pi^2/12}.
        \]

  \item % 1.6.30
        Misalkan $(a_n)_{n \ge 1}$ adalah barisan aritmetika dengan suku-suku bilangan positif. Hitunglah
        \[
          \lim_{n \to \infty} \frac{n(a_1 \dots a_n)^{1/n}}{a_1 + \dots + a_n}.
        \]
\end{enumerate}

\section*{Aljabar Linier}
\begin{enumerate}[label=\arabic*.]
  \item % Problem 1
        Misalkan $X$ adalah matriks tak singular dengan kolom-kolom $X_1, X_2, \ldots, X_n$. Misalkan $Y$ adalah matriks dengan kolom-kolom $X_2, X_3, \ldots, X_n, 0$. Tunjukkan bahwa matriks $A = YX^{-1}$ dan $B = X^{-1}Y$ memiliki rank $n - 1$ dan hanya memiliki nol untuk nilai eigennya.

  \item % Problem 3
        Misalkan $A$ dan $B$ adalah matriks real $n \times n$ sedemikian sehingga $A^2 + B^2 = AB$. Buktikan bahwa jika $BA - AB$ adalah matriks yang dapat dibalik (invertibel) maka $n$ habis dibagi 3.
\end{enumerate}

\section*{Analisis Kompleks}
\begin{enumerate}[label=\arabic*.]
  \item % Problem 14
        Misalkan $x, y, z$ adalah bilangan kompleks.
        \begin{enumerate}[label=(\alph*)]
          \item Buktikan bahwa

                \[
                  |x| + |y| + |z| \le |x+y-z| + |x-y+z| + |-x+y+z|.
                \]
          \item Jika $x, y, z$ berbeda dan berturut-turut $x+y-z, x-y+z, -x+y+z$ memiliki nilai mutlak yang sama, buktikan bahwa
                \[
                  2(|x| + |y| + |z|) \le |x+y-z| + |x-y+z| + |-x+y+z|.
                \]
        \end{enumerate}

  \item % Problem 5
        Selesaikan persamaan
        \[
          \cos x + \cos 2x - \cos 3x = 1.
        \]
\end{enumerate}

\section*{Kombinatorika}
\begin{enumerate}[label=\arabic*.]
  \item % Meja Bundar Bubur Ayam
        Sebanyak 8 orang duduk di meja makan yang berbentuk lingkaran. Setiap orang dihidangkan dengan sebuah kotak makanan yang berisi bubur ayam atau lontong sayur. Apabila peluang bahwa tidak ada dua orang bersebelahan yang mendapatkan bubur ayam dapat dinyatakan dalam bentuk $\frac{m}{n}$ untuk bilangan asli relatif prima $m$ dan $n$, tentukan nilai dari $1000m + n$.

  \item % Bertukar Kado
        Terdapat 6 anak, A, B, C, D, E, dan F, akan saling bertukar kado. Tidak ada yang menerima kadonya sendiri, dan kado dari A diberikan kepada B. Banyaknya cara membagikan kado dengan cara demikian adalah \dots

  \item % Seratus Bilangan Bulat
        Seratus bilangan bulat disusun mengelilingi lingkaran sedemikian sehingga (menurut arah jarum jam) setiap bilangan lebih besar daripada hasil penjumlahan dua bilangan sebelumnya. Maksimal banyaknya bilangan bulat positif yang terdapat pada lingkaran tersebut adalah \dots

  \item % Deret Kombinasi
        Untuk setiap bilangan asli $n \in \mathbb{N}$ dengan $n \ge 2$, nilai dari
        \[
          \frac{1}{n} \binom{n}{1} + \frac{2}{n} \binom{n}{2} + \frac{3}{n} \binom{n}{3} + \dots + \frac{n-1}{n} \binom{n}{n-1}
        \]
        adalah \dots
\end{enumerate}

\section*{Struktur Aljabar}
\begin{enumerate}[label=\arabic*., start=37]
  \setcounter{enumi}{0} % Agar mulai dari 1
  \item % 37
        Diberikan $\alpha : G \to G_1$ dan $\beta : G \to G_2$ adalah homomorfisma grup.
        \begin{enumerate}[label=(\alph*)]
          \item Tunjukkan bahwa jika $\ker \alpha \subseteq \ker \beta$ dan $\alpha$ surjektif, maka terdapat homomorfisma grup $\phi : G_1 \to G_2$ yang terdefinisi dengan baik sedemikian sehingga $\beta = \phi \circ \alpha$.
          \item Tunjukkan bahwa jika $\ker \alpha \not\subseteq \ker \beta$, maka tidak ada homomorfisma $\phi$ yang memenuhi kondisi tersebut.
        \end{enumerate}

  \item % 38
        Buktikan atau berikan contoh penyangkal bahwa untuk setiap $\sigma \in A_5$, terdapat $\tau \in S_5$ sedemikian sehingga $\tau^2 = \sigma$.

  \item
        Dapatkan semua bilangan prima $p$ sehingga persamaan
        \begin{align*}
          \sigma^p & = (1 \quad 2 \quad \dots \quad p)(p+1 \quad p+2 \quad \dots \quad 2p)                                                                              \\
                   & = \begin{pmatrix} 1 & 2 & \dots & p & p+1 & p+2 & \dots & 2p \\ 2 & 3 & \dots & 1 & p+2 & p+3 & \dots & p+1 \end{pmatrix}, \quad \sigma \in S_{2p}
        \end{align*}
        memiliki solusi. Dapatkan semua solusinya untuk setiap bilangan prima $p$ yang memenuhi.

  \item
        Suatu subgrup $H$ dari $G$ disebut subgrup \textit{normal} jika untuk setiap $g \in G$ dan $h \in H$ berlaku $g^{-1}hg \in H$. Suatu subgrup $H$ dari $G$ dikatakan mempunyai sifat \textit{refleksif} jika untuk setiap $a,b \in G$ dengan $ab \in H$ berlaku $ba \in H$. Tunjukkan bahwa $H$ merupakan subgrup normal jika dan hanya jika $H$ bersifat refleksif.

  \item % 59
        Jika $G$ adalah grup siklik dengan tak hingga unsur dan $H$ subgrup tak trivial dari $G$, tunjukkan bahwa $[G : H]$ berhingga.

  \item % 60
        Misalkan $G$ grup komutatif dan $g \in G$ dengan $|g| = 2$. Jika $H$ subgrup dari $G$, buktikan bahwa $H \cup gH$ juga subgrup $G$. Bagaimana jika $G$ tak komutatif?

  \item % 61
        Misalkan $H$ subgrup dari suatu grup $G$ dengan $[G : H] = 2$. Jika $a,b \in G - H$, buktikan bahwa $ab \in H$.

  \item % 62
        Misalkan $R$ adalah suatu ring. Unsur $a \in R$ disebut \textbf{nilpoten} jika ada bilangan asli $n$ sehingga $a^n = 0$.
        \begin{enumerate}[label=(\alph*)]
          \item Tunjukkan bahwa jika $a, b$ nilpoten di $R$ dan $ab = ba$, maka $a+b$ juga nilpoten. Bagaimana jika $ab \neq ba$?
          \item Tunjukkan bahwa $m \in Z_n$ nilpoten jika dan hanya jika $m$ habis dibagi setiap faktor prima dari $n$.
        \end{enumerate}
\end{enumerate}

\end{document}